%-----------------------------------
\section{Pesquisa experimental}
\label{sec_tipo_experimental}

\subsection{Caracterização}

A pesquisa experimental manipula deliberadamente uma ou mais variáveis (variáveis independentes ou fatores) e mede o efeito sobre outras (variáveis dependentes), mantendo as demais condições sob controle \cite{wohlin2012experimentation}. Em Sistemas de Informação, são exemplos a comparação de desempenho entre bancos de dados ou \textit{frameworks}, a avaliação de algoritmos de aprendizado de máquina \cite{Lecun2015DeepLearning} e experimentos controlados com usuários.

O ponto central é a \emph{reprodutibilidade}: outro pesquisador deve ser capaz de repetir o experimento a partir da descrição apresentada. Por isso, ambiente, dados, parâmetros e número de repetições devem ser informados.

\subsection{Estrutura típica do trabalho}

\begin{quadro}[htb]
	\centering
	\caption{Estrutura típica de um TCC de pesquisa experimental.}
	\label{qua_estrutura_experimental}
	\small
	\begin{tabularx}{\textwidth}{|l|X|}
		\hline
		\textbf{Capítulo} & \textbf{Seções típicas} \\ \hline
		Metodologia & Questões de pesquisa e hipóteses; objetos de estudo ou conjuntos de dados; variáveis e métricas; protocolo experimental (ambiente, repetições); análise estatística. \\ \hline
		Resultados e discussão & Resultados por questão de pesquisa; teste das hipóteses; discussão; ameaças à validade. \\ \hline
		Apêndices & Configuração do ambiente; \textit{scripts} e código-fonte (ou link para repositório); dados brutos. \\ \hline
	\end{tabularx}
	\fonte{Elaborado pelo autor.}
\end{quadro}

O Quadro~\ref{qua_estrutura_experimental} resume a estrutura desse tipo de trabalho. As hipóteses costumam ser formuladas em pares, com uma hipótese nula ($H_0$) e uma alternativa ($H_1$).

\subsection{Exemplo: comparação de desempenho}

Considere um experimento que compara o tempo de resposta de dois bancos de dados. Para cada um, a consulta é executada $n$ vezes e calcula-se a média amostral pela Equação~\eqref{eq_media} e o desvio padrão amostral pela Equação~\eqref{eq_desvio}.

\begin{equation}
	\bar{x} = \frac{1}{n} \sum_{i=1}^{n} x_i
	\label{eq_media}
\end{equation}

\begin{equation}
	s = \sqrt{\frac{1}{n-1} \sum_{i=1}^{n} \left( x_i - \bar{x} \right)^2}
	\label{eq_desvio}
\end{equation}

O procedimento de medição pode ser descrito em pseudocódigo, como no Algoritmo~\ref{alg_medicao}, e os resultados, apresentados em tabela, como a Tabela~\ref{tab_resultados_experimento}.

\begin{algorithm}[H]
	\SetAlgoLined
	\KwIn{banco de dados $B$, consulta $q$, número de repetições $n$}
	\KwOut{lista de tempos $T$}
	$T \leftarrow \emptyset$\;
	executar $q$ em $B$ uma vez para aquecimento do \textit{cache}\;
	\For{$i \leftarrow 1$ \KwTo $n$}{
		$t_0 \leftarrow$ tempo atual\;
		executar $q$ em $B$\;
		adicionar (tempo atual $- t_0$) a $T$\;
	}
	\Return{$T$}\;
	\caption{Medição do tempo de resposta}
	\label{alg_medicao}
\end{algorithm}

\begin{table}[htb]
	\centering
	\caption{Tempo de resposta por banco de dados, em milissegundos (dados fictícios).}
	\label{tab_resultados_experimento}
	\begin{tabular}{lrrrr}
		\toprule
		Banco de dados & $n$ & $\bar{x}$ & $s$ & Mediana \\
		\midrule
		Banco A & 30 & 12,40 & 1,85 & 12,10 \\
		Banco B & 30 & 18,75 & 3,20 & 18,05 \\
		\bottomrule
	\end{tabular}
	\fonte{Elaborado pelo autor.}
\end{table}
